% TOP_PROBLEM: {"id": 6200006, "title": "Boundaries of Groups and Kleinian Groups — Problem 6", "category_id": 6, "category": {"id": 6, "name": "geometry", "display_name": "Geometry", "description": "Euclidean and non-Euclidean geometry, geometric structures.", "slug": "geometry", "order_index": 6, "created_at": "2026-07-31T15:26:25.671Z"}, "status": "open", "problem_number": "AMR-061-0006", "set_id": 13}
% FIRST_POSTED: 2026-10-02
% ENTRY_KIND: statement_only
% UnsolvedMath ID 6200006; source catalogue code AMR-061-0006
% !TeX program = lualatex
% Definition and sources only; this file is not an LLM solution attempt.
\documentclass[11pt,a4paper]{article}
\usepackage[margin=25mm]{geometry}
\usepackage{fontspec}
\setmainfont{lmroman10-regular.otf}[BoldFont=lmroman10-bold.otf,
 ItalicFont=lmroman10-italic.otf,BoldItalicFont=lmroman10-bolditalic.otf]
\usepackage{amsmath,amssymb,mathtools}
\usepackage{unicode-math}
\setmathfont{latinmodern-math.otf}
\AtBeginDocument{\let\setminus\smallsetminus}
\usepackage{xurl,hyperref}
\hypersetup{colorlinks=true,urlcolor=blue}
\setcounter{secnumdepth}{0}
\setlength{\parindent}{0pt}
\setlength{\parskip}{5pt}
\setlength{\emergencystretch}{3em}
\begin{document}
\section*{Boundaries of Groups and Kleinian Groups — Problem 6}\label{6200006}
{\small UnsolvedMath ID 6200006 · AMR-061-0006 · Geometry · AMR Open Problem Lists}
\subsection{Definitions and mathematical statement}
A finitely generated group $G$ is hyperbolic relative to a finite family $\mathcal P=\{P_1,\ldots,P_s\}$ in the Bowditch sense if it acts on a connected fine hyperbolic graph with finitely many edge orbits and finite edge stabilizers, whose infinite vertex stabilizers are exactly conjugates of the $P_i$. Fine means that each edge belongs to only finitely many embedded circuits of any specified length; hyperbolic means that geodesic triangles are uniformly thin.

The Bowditch boundary is the compact boundary obtained from the graph's Gromov boundary by adjoining its infinite-valence parabolic vertices, with the usual observer topology. Its equivariant homeomorphism type is independent of the chosen graph for this relatively hyperbolic structure.

Does there exist such a group with nonempty peripheral family, each $P_i$ finitely generated nilpotent of class at least three, and
\[
 \partial_B(G,\mathcal P)\cong S^n
\]
for some $n\geq1$? Nilpotency class is defined using the lower central series $\gamma_1(P)=P$, $\gamma_{j+1}(P)=[P,\gamma_j(P)]$: class $c$ means $\gamma_{c+1}(P)=1$ but $\gamma_c(P)\ne1$.

The sphere is a topological sphere; no round metric or conformal action is required. The nilpotent groups must occur as peripheral subgroups of the given structure, rather than merely as arbitrary subgroups of $G$. Abelian and two-step nilpotent cusp groups do not meet the requested class threshold.
\subsection{Short English statement}
Can a spherical relatively hyperbolic boundary coexist with nilpotent peripheral groups of class at least three?
\subsection{Sources}
\begin{itemize}
\item[S1] ULAM.AI and the UnsolvedMath Contributors, supplied problems.json, record 6200006 (AMR-061-0006), AMR Open Problem Lists. Catalogue identity and source transcription; CC BY 4.0.
\url{https://huggingface.co/datasets/ulamai/UnsolvedMath}.
\item[S2] Kapovich - Problems on Boundaries of Groups and Kleinian Groups (2005), Problem 6, PDF page 3. Original source indexed by AMR-061-0006.
\url{https://www.math.ucdavis.edu/~kapovich/EPR/problems.pdf}.
\end{itemize}
\end{document}
